\chapter{归并与动态输入数据识别}

\begin{solutionexercise}{1}
\problemheading 假设要归并列表
$A=\{1,7,8,9,10\}$ 与 $B=\{7,10,10,12\}$。$C[8]$ 的两个协同秩是多少？

\approachheading 协同秩 $(i,j)$ 满足 $i+j=k$，且长度为 $k$ 的输出前缀恰由
$A[0..i-1]$ 与 $B[0..j-1]$ 稳定归并而成；相等时令 $A$ 先于 $B$。

\answerheading 完整稳定归并为
\[
[1,7_A,7_B,8,9,10_A,10_{B_1},10_{B_2},12].
\]
前 8 个元素已经耗尽 $A$ 的 5 个元素以及 $B$ 的前 3 个元素，故
\[
 i=5,\qquad j=8-i=3.
\]
边界条件也成立：$A[i-1]=10\le B[j]=12$；$i=m$ 使第二个条件自动成立。

\complexity{二分协同秩为 $O(\log(\min(m,n)+1))$，手算为常数}{$O(1)$}{$i+j=k$ 的候选可二分搜索}
\conclusionheading 两个协同秩为 $\boxed{(5,3)}$。
\discussionheading
\discussionpoint{协同秩把输出位置变成输入切分。}若归并结果的前 $k$ 项消费了 $A$ 的 $i$ 项和 $B$ 的 $j$ 项，就必有 $i+j=k$，所以所有候选都落在归并矩阵的同一条对角线上。合法点还要满足切分左侧元素不大于下一侧候选，保证两段输入前缀合在一起恰是最小的 $k$ 项。找到这个点后，输出区间 $[k_0,k_1)$ 只需读取两端协同秩之间的输入片段，可以与其他区间独立归并。例如本题 $k=8$ 的 $(5,3)$ 表示无论两数组长度如何失衡，前八个输出的消费总数始终精确为八。它把按输入切块造成的数据相关负载不均，转化成按输出数量均匀分工，也是 merge path 几何图像的核心。

\discussionpoint{稳定性隐藏在严格与非严格比较的方向中。}当两个键相等并约定 $A$ 中记录先于 $B$，边界的一侧必须允许等号，另一侧必须严格，否则同一对角线上可能出现两个合法候选或一个也没有。以 $A=[2_A]$、$B=[2_B]$、$k=1$ 为例，稳定切分应消费 $A$ 而不是 $B$；只看输出键值都是 2，错误完全不可见。给记录附上来源与原始下标后，比较规则才成为可观察契约。也可把比较键形式化为 $(key,sourcePriority,index)$，使手算、CPU 参考和 GPU 边界判断共享同一个全序。相同的 tie-break 还决定数据库稳定连接、时间序列同刻事件和外部排序多路归并的顺序，因此不是实现细节。

\discussionpoint{概念哨兵不能变成越界解引用。}数学证明常把数组左外侧写作 $-\infty$、右外侧写作 $+\infty$，代码却应先判断 $i=0,m$ 或 $j=0,n$，再决定是否访问邻居，不能真的形成并解引用越界地址。搜索区间至多包含较短输入长度再加一个边界候选，所以复杂度写为 $O(\!\log(\min(m,n)+1))$ 才在某侧为空时仍有定义；直接写 $O(\log\min(m,n))$ 会出现 $\log0$。当 tile 只有几项，线性推进或复用上一个切分反而比二分的分支更便宜，说明渐近最优不等于每个局部规模都最优。

\discussionpoint{不变量测试能覆盖归并路径的整个应用族。}每个切分都应同时满足 $i+j=k$、下标范围和两条有序边界，所有相邻区间拼接后还必须无重无漏地等于稳定串行归并。测试应包含空数组、长度悬殊、完全分离和大量重复键，因为它们分别触发对角线端点、窄搜索区间和 tie-break。特别是某侧为空时，搜索区间仍有一个合法边界候选，复杂度应写成 $O(\log(\min(m,n)+1))$，不能让公式出现 $\log0$。相同切分方法可用于数据库 sort-merge join、倒排表合并和有序流重分区，但记录负载可能比键大得多。将切分正确性与实际搬运字节分开验证，才能既保证排序语义又评价负载均衡。
\variantheading 若改求 $k=6$ 的边界，前六项为 $[1,7_A,7_B,8,9,10_A]$，故协同秩为
$\boxed{(5,1)}$；这与本题 $k=8$ 的 $(5,3)$ 不同。
\practiceheading 可用 Thrust \code{stable_merge} 或 CPU 稳定归并生成黄金输出，并随机加入
大量重复键检查每个 $k$ 的 $i+j=k$、左右边界条件及相邻分区无重叠。
\sourcecheck{https://doi.org/10.1007/978-3-642-33518-1_25}{Siebert 与 Träff 的稳定并行归并原始论文；高置信度}
假设审查：$C[8]$ 表示“秩为 8 的边界”，协同秩描述它之前的 8 项，并非声称 $C[8]$
本身来自 $A[5]$ 或 $B[3]$。若相等时改让 $B$ 优先，中间三个 10 的来源顺序会变，但该
边界仍是 $(5,3)$。
\end{solutionexercise}

\begin{solutionexercise}{2}
\problemheading 完成原书图 13.6 中线程 2 的协同秩函数计算过程。为使题目可独立理解，
该图归并的输入为 $A=[1,7,8,9,11]$、$B=[7,10,10,12]$，并用 3 个线程、每线程
负责 3 个连续输出。线程 2 负责 $C[6]$ 至 $C[8]$，因此以
$(k,A,m,B,n)=(6,A,5,B,4)$ 调用协同秩函数；请给出二分搜索各轮的候选
$(i,j)$、下界以及最终协同秩。

\approachheading 严格执行书中 \code{co_rank(k,A,m,B,n)}：初始
$i=\min(k,m)$、$j=k-i$，并保持 $i+j=k$。

\answerheading $m=5,n=4,k=6$，初值与下界为
\[
(i,j)=(5,1),\qquad i_{low}=\max(0,6-4)=2,\quad
j_{low}=\max(0,6-5)=1.
\]
第 0 轮：$A[4]=11>B[1]=10$，故
$\delta=\lceil(5-2)/2\rceil=2$，更新为
$(i,j)=(3,3)$、$j_{low}=1$。

第 1 轮：$A[2]=8\le B[3]=12$，但 $B[2]=10\ge A[3]=9$，故
$\delta=\lceil(3-1)/2\rceil=1$，更新为
$(i,j)=(4,2)$、$i_{low}=3$。

第 2 轮：$A[3]=9\le B[2]=10$ 且 $B[1]=10<A[4]=11$，两个停止条件均满足，返回
$i=4$，从而 $j=2$。前 6 项确为 $[1,7_A,7_B,8,9,10_{B_1}]$。

\editorialnote 英文底本和中文正文紧邻该图的一段叙述把 $k=6$ 的结果写成 $(5,1)$，并
声称线程 1 输出 $[8,9,11]$；其后线程才输出 10，会使 $C$ 不是有序数组。按底本给出的
函数逐行执行得到 $(4,2)$，上面的反例也独立证实本题应采用这一结果。

\complexity{$O(\log(\min(m,n)+1))$ 时间}{$O(1)$}{不同线程可独立求各自边界}
\conclusionheading 线程 2 的正确协同秩为 $\boxed{(4,2)}$，共经历三次条件检查轮。
\discussionheading
\discussionpoint{先用不变量审图，而不是让图替算法作证。}一条可信协同秩轨迹中的每个候选都必须在目标对角线上，即 $i+j=k$；终点还要满足稳定归并的左右夹逼条件。若教材箭头指到不满足这些式子的格点，应最小修正图示，而不能为了匹配图而改动算法边界。比如目标 $k=6$ 时，图中任何 $(4,3)$ 都因和为 7 而立即失效，无需继续争论比较结果。只从最终合并序列倒推可能掩盖中间搜索区间曾扩大、比较方向曾写反等问题。逐轮核对不变量能指出究竟是候选坐标、上下界还是严格符号出错，这也是审查论文伪代码和硬件波形的通用思路。

\discussionpoint{重复键是观察隐藏稳定约定的探针。}若两个输入都含键 5，选择 $5_A$ 或 $5_B$ 作为第一个输出在数值数组中完全一样，但会给出不同协同秩，并影响携带载荷的记录次序。测试数据应把元素写成 $(key,source,index)$，期望次序便能明确规定同键时先比较来源或原始下标。严格与非严格不等式必须与这条规则成对设计，不能两侧都严格或都宽松。若载荷分别是“旧版本”和“新版本”，一次隐藏换序甚至会改变下游去重保留哪条记录，而不只是展示顺序。数据库同键事件、日志时间戳和稳定排序都依赖相同约定；通过来源标签验证后，稳定性才从口头承诺变成可执行性质。

\discussionpoint{二分轨迹是压缩而高价值的调试证据。}每轮只需记录搜索下界、上界、候选 $(i,j)$ 以及哪条边界比较失败，就能检查区间是否严格缩小、候选是否始终满足 $i+j=k$。若循环卡住，轨迹会直接显示整数中点取整后上下界未变化；若终点越界，则可追到第一次非法更新。合法实现每轮至少排除一个候选，因此迭代上界可由初始候选数的对数给出，超过上界本身就是错误信号。GPU 上不应让成千线程同时 \code{printf}，可只让失败块领导者原子保留一个紧凑样例，随后在 CPU 重放完整搜索。类似的“设备捕获最小反例、主机确定性复现”也适合哈希探测和并行图算法调试。

\discussionpoint{勘误记录本身应当可复现。}修正底本时应保留原图号、原坐标或比较符号、修正结果和依据，使读者能从同一输入重新走出差异；静默改图会切断证据链。最小复现应同时保存两输入、目标 $k$、稳定规则和逐轮区间，因为只截终点无法区分图错与算法错。生产代码也不能只保存这一个示例，因为固定数据可能恰好没有重复键或端点。应把对角线、范围和夹逼条件写成调试断言，再用属性测试生成空侧、重复键和极端长度比。这样一次教材勘误可以转化为长期回归规则，未来更换索引宽度、tile 大小或比较器时仍能捕获同类错误。
\variantheading 若把 $A[4]$ 从 11 改为 10，$k=6$ 的稳定前缀仍先取 $A$ 的三个 10 之前
可用项，协同秩变为 $(5,1)$；停止条件必须继续让 $A$ 的相等键优先。
\practiceheading 把底本示例和随机重复键数据加入回归测试，逐个输出边界与串行稳定归并
对比；GPU 上再用 Compute Sanitizer memcheck 检查 $i=0,m$ 或 $j=0,n$ 的哨兵分支。
\sourcecheck{https://davidbader.net/publication/2012-gm-ba/2012-gm-ba.pdf}{GPU Merge Path 作者公开版对对角划分/协同秩思想的独立佐证；并以 DOI 10.1145/2304576.2304621 核对正式出版信息；高置信度}
对抗审查：若用不严格条件 $B[j-1]\le A[i]$，相等键可能同时被相邻分区认领，破坏稳定性；
书中第二条件用 $\ge$ 触发调整正是为了让 $A$ 的相等键优先。
\end{solutionexercise}

\begin{solutionexercise}{3}
\problemheading 对原书图 13.12 中载入 $A$、$B$ 输入块的两个 for 循环增加协同秩
函数调用，使其只载入当前 while 循环迭代实际会消费的 $A$、$B$ 元素。原内核每轮最多
生成 \code{tile_size} 个输出，并分别从尚未消费的两个有序输入中各载入至多
\code{tile_size} 项；需要先由本轮输出边界的协同秩确定两个输入的精确消费量，再由块内
线程协作载入，同时正确处理最后一轮和任一输入先耗尽的边界情形。

\approachheading 当前轮最多生成 \code{tile_size} 个输出。先由一个线程对尚未消费的两个
全局子数组求这个输出边界的协同秩，得到精确的 \code{A_take/B_take}，再合作加载。

\answerheading 以下代码替换原书图 13.12 的两段固定上限循环，并相应更新消费计数：
\begin{lstlisting}[language=C++]
__shared__ int A_take_s, B_take_s, C_take_s;

while (C_completed < C_length) {
  if (threadIdx.x == 0) {
    int remainC = C_length - C_completed;
    C_take_s = min(tile_size, remainC);
    int remainA = A_length - A_consumed;
    int remainB = B_length - B_consumed;
    A_take_s = co_rank(C_take_s,
        A + A_curr + A_consumed, remainA,
        B + B_curr + B_consumed, remainB);
    B_take_s = C_take_s - A_take_s;
  }
  __syncthreads();

  for (int p=threadIdx.x; p<A_take_s; p+=blockDim.x)
    A_S[p] = A[A_curr + A_consumed + p];
  for (int p=threadIdx.x; p<B_take_s; p+=blockDim.x)
    B_S[p] = B[B_curr + B_consumed + p];
  __syncthreads();

  int c0 = int((static_cast<unsigned long long>(threadIdx.x) * C_take_s)
               / blockDim.x);
  int c1 = int((static_cast<unsigned long long>(threadIdx.x + 1) * C_take_s)
               / blockDim.x);
  int a0 = co_rank(c0,A_S,A_take_s,B_S,B_take_s), b0=c0-a0;
  int a1 = co_rank(c1,A_S,A_take_s,B_S,B_take_s), b1=c1-a1;
  merge_sequential(A_S+a0,a1-a0,B_S+b0,b1-b0,
                   C+C_curr+C_completed+c0);
  __syncthreads();
  A_consumed += A_take_s; B_consumed += B_take_s;
  C_completed += C_take_s;
  __syncthreads();
}
\end{lstlisting}
协同秩定义保证当前输出前缀恰好消费 \code{A_take+B_take=C_take} 个输入；共享内存
中的线程级归并再把这一前缀无重叠地划分。循环末屏障防止快线程覆盖仍在使用的 tile。

\complexity{$O(m+n)$ 归并工作，每轮另有 $O(\log(m+n))$ 边界搜索}{$O(\text{tile\_size})$ 共享空间}{块间、块内线程区间均可并行}
\conclusionheading 全局加载总量从每轮“至多两个完整 tile”降为恰好当前轮消费的元素数。
\discussionheading
\discussionpoint{输出预算可以精确反推出本轮输入需求。}若共享 tile 本轮只准备产生 $C_{take}$ 个结果，从当前输入游标出发求终止对角线的协同秩，就得到恰好要消费的 $A_{take}$ 与 $B_{take}$，并满足两者之和等于输出数。朴素方案若总从两侧各加载一个完整 tile，在两个数组完全分离或一侧接近结束时，会搬入大量本轮不会比较的元素。按协同秩裁剪后，加载集合与消费集合一致，节省的是全局字节而非归并比较本身。相同“从输出容量反推输入边界”的思想还用于分页 join、流窗口和压缩块解码。

\discussionpoint{一次二分搜索只有在省下的通信更贵时才值得。}领导线程执行 $O(\!\log T)$ 比较，还要把消费数广播给全块并同步；若 tile 很小、数据已命中缓存或两侧消费恰好均衡，控制成本可能超过少量多余加载。tile 较大、全局内存昂贵或输入高度偏斜时，避免半个 tile 的无用传输则容易覆盖搜索成本。可用 $T_{saved}\approx bytes_{avoided}/B_{eff}$ 与二分、广播和屏障时间比较盈亏点，并复用前一轮切分作为下一轮搜索窄区间。这个模型把算法优化从“少加载一定快”变成可在目标设备上验证的条件判断。

\discussionpoint{协作加载必须覆盖任意消费数并维护流水生命周期。}让线程 $t$ 处理 $t,t+B,t+2B,\ldots$ 的局部下标，可覆盖从 0 到任意 \code{count-1} 的全部元素，即使数量不是块大小倍数或最后一轮只有几项，也不会遗漏。例如 $B=128,\code{count}=130$ 时，线程 0 和 1 各再搬一项，其余线程只搬一项，循环无需特殊尾部分支。支持异步全局到共享拷贝时，可用双缓冲让下一轮加载与当前轮归并重叠；但覆盖旧缓冲前必须确认所有消费者完成，提交组和等待点也要与该轮边界元数据配对。若只同步数据不更新消费数，线程可能用新 tile 配旧边界。这个版本化问题与网络环形缓冲和生产者--消费者队列完全相似。

\discussionpoint{输入分布决定裁剪加载的实际收益。}均匀交错时两侧消费大致均衡，完全分离时一轮可能几乎只取一侧，全相等又会把比例交给稳定 tie-break，一侧很短则频繁触发端点；四类数据覆盖了主要路径。若每轮输出 $T$ 项而朴素地两侧各载入 $T$ 项，完全分离时裁剪最多可避免近 $T$ 项无用搬运，均匀交错时收益则小得多。性能报告应同时给出实际加载元素数、全局扇区、二分迭代、屏障等待和结果吞吐，才能确认少搬元素确实转化为少事务。正确性还要累计所有轮的 $A_{take}$、$B_{take}$，最终分别等于输入长度，并与稳定 CPU 归并逐项相同。这样数据分布既是性能变量，也是边界证明的测试生成器。
\variantheading 若当前轮生成 1000 项、块含 128 个线程，则线程 127 的区间为
\[
[\lfloor127\times1000/128\rfloor,1000)=[992,1000),
\]
所以最后 8 项不会因整除截断而丢失。
\practiceheading 用 CPU 的标准归并对照不同长度、全相等键和空侧输入，并重点覆盖 tile
大小不能整除块大小的组合；Nsight Compute 可比较修复前后的实际加载字节与领导
线程二分开销。
\sourcecheck{https://davidbader.net/publication/2012-gm-ba/2012-gm-ba.pdf}{GPU Merge Path 作者公开版的分区与分块归并方法；并以 DOI 10.1145/2304576.2304621 核对正式出版信息；高置信度}
边界审查：末轮 \code{C_take} 可能小于块大小，必须用它限制线程输出；若只缩短加载却仍
把 \code{tile_size} 传给协同秩，会读取未初始化共享内存。
\end{solutionexercise}

\begin{solutionexercise}{4}
\problemheading 考虑并行归并两个长度分别为 1,030,400 和 608,000 的数组。假设每线程
归并 8 个元素，线程块大小为 1024。分别回答：
\begin{enumerate}[label=(\alph*)]
  \item 在原书图 13.9 的基本归并内核中，有多少线程在全局内存数据上执行二分搜索？
  \item 在原书图 13.11 至图 13.13 的分块归并内核中，有多少线程在全局内存数据上执行
  二分搜索？
  \item 在同一分块归并内核中，有多少线程在共享内存数据上执行二分搜索？
\end{enumerate}

\approachheading 先求输出总长、线程数和块数，再区分“执行搜索的不同线程数”与“函数
调用次数”。

\answerheading
\[
N=1{,}030{,}400+608{,}000=1{,}638{,}400,
\]
\[
T=N/8=204{,}800,\qquad G=T/1024=200.
\]

基本内核中每个线程都为输出区间两端调用协同秩，所以在全局内存数据上搜索的是
$\boxed{204{,}800}$ 个不同线程（共 $409{,}600$ 次调用）。

分块内核只让每块一个领导线程求块级起止协同秩，所以在全局数据上搜索的是
$\boxed{200}$ 个不同线程（共 400 次调用）。块内每个线程仍需在共享 tile 上求自己的
起止边界，所以在共享数据上搜索的是 $\boxed{204{,}800}$ 个不同线程（共 409,600 次
调用；本题参数使每线程的 8 项恰在一个块级轮次中完成）。

\complexity{搜索本身每次 $O(\log N)$；随后归并总工作 $O(N)$}{$O(\text{tile size})$ 共享空间/块}{$204800$ 个线程、200 个块}
\conclusionheading (a) 204,800；(b) 200；(c) 204,800。
\discussionheading
\discussionpoint{两级切分把昂贵搜索提升到块级共享。}块领导者先在完整全局数组上求当前输出 tile 的两端协同秩，确定要搬入共享内存的输入片段；块内线程再在这两个小片段上为自己的输出子区间求局部切分。于是高延迟、可能跨许多缓存行的全局二分从每线程一次降为每块一次，增加的是低延迟共享搜索。若 1024 个线程各自做全局二分，搜索起点就有 1024 份；提升后每块只需求少数 tile 边界，再让线程在共享窗口内细分。结构上它像页表和缓存的两级定位：粗层缩小工作集，细层在近处完成个体分工。正确性来自局部坐标加回全局起点后仍满足同一对角线与夹逼不变量。

\discussionpoint{线程数、调用数和比较数必须分开统计。}一个线程可能为左右两个边界调用两次二分，一次二分又可能执行若干比较；“有多少线程搜索”不能直接回答总比较工作。尾块仍会启动完整线程块，但超出实际输出范围的线程可能跳过搜索或只处理空区间，取决于边界代码。若每线程生成多个输出，线程数还会少于输出数。粗略模型可写成 $C\approx Q\bar L+C_{block}$，其中 $Q$ 是有效调用数，$\bar L$ 是每次平均迭代，而不是把启动线程数当作比较数。这个计数纪律同样适用于哈希探测、BVH 遍历和数据库索引查询。

\discussionpoint{tile 尺寸在搜索频率与片上驻留之间取平衡。}较小 tile 需要更多块级全局切分和同步，输入片段复用也有限；较大 tile 减少边界搜索，却消耗更多共享内存，并可能让每 SM 只能驻留很少块。若 tile 从 1024 翻倍到 2048 项，边界搜索数近似减半，但双缓冲共享容量也近似翻倍，驻留块数可能发生阶跃下降。记录若包含大 payload，搬入整个记录会进一步放大容量，可只缓存键和索引，归并确定后再搬 payload。最优点还受比较器成本和输入长度比影响，不能从整数键实验直接外推。通过扫描 tile 大小并记录全局搜索、共享占用和端到端吞吐，才能找到真正的资源拐点。

\discussionpoint{分层归并连接到存储系统中的有序数据维护。}数据库 sort-merge join 用有序键对齐记录，LSM compaction 合并不同层的有序段，倒排索引交集处理长度高度不均的 posting list；它们都可用对角切分分配输出工作。实际记录比较可能访问字符串，payload 搬运也远大于整数，所以键搜索、记录物化与写带宽应分别计量。尤其 LSM 压实常受写放大限制，即使切分让线程工作均衡，额外物化一遍 payload 仍可能吞掉全部收益。基准要改变输入长度比、重复键和 payload 大小，并用稳定 CPU 结果核对来源次序。这样才能判断两级切分究竟缓解了负载不均，还是被比较与搬运成本掩盖。
\variantheading 若每线程改为归并 16 项而块仍为 1024 线程，则共有 102,400 个线程和
100 个块；全局搜索线程数为 100，共享搜索线程数为 102,400。
\practiceheading 在生产基准中分别统计全局与共享内存访问、二分分支效率和块占用率，并用
GPU Merge Path 或 Thrust merge 作基线；长度乘加应使用 \code{size_t}，避免大输入的 32 位溢出。
\sourcecheck{https://davidbader.net/publication/2012-gm-ba/2012-gm-ba.pdf}{GPU Merge Path 作者公开版对层次化分区搜索的分析；并以 DOI 10.1145/2304576.2304621 核对正式出版信息；高置信度}
假设审查：题目问“多少线程”而不是“多少次搜索”。若实现复用相邻线程共享的区间端点，
调用次数可以再减；这里严格按原书两个内核的源代码计数。
\end{solutionexercise}
